\scanpage{85}
\mspage{84}
C'est cette intuition que je voudrais maintenant essayer de cerner de plus près.

\figure{\begin{tikzpicture}\draw (0,.5)--(.25,0);\end{tikzpicture}}

Soit $\unsure{\mathscr{F}|\mathscr{E}}$ un atlas, $F$ une figure, $\widetilde F=I$ l'ens ordonné de ses strates.

\deleted{Dans le cas où $\operatorname{supp}_{\mathscr{F}}F=\unsure{J|I}$}

Intuitivement, on va identifier les éléments de $I$ à des ``strates ouvertes'', et on veut identifier $\unsure{\mathcal{P}|\mathscr{P}}(I)$ à \deleted{quelque} un ensemble de supports (réunions de ``strates ouvertes'').

Soit $X$ une strate. On a envie de donner un sens intrinsèque à ``$X\smallsetminus\unsure{\partial|\delta}X$'', L'idée est de \deleted{prendre} définir
\[
X^{\circ}=\operatorname{supp}_{\unsure{M|\mu}}X\cap\operatorname{cosupp}_M\partial X\qquad(\subset M)
\]
où je rappelle
\[
\partial X=\{Y\in M\mid Y\leftarrow X\}=\widetilde X\smallsetminus\{X\}
\]
On aimerait déjà vérifier que
\[
\operatorname{supp}X=\underset{Y\in X}{\operatorname{Sup}}\,Y^{\circ}
\]
\margin{Sup dans $\overline{\Sigma}_M$}
i.e, passant aux supports complémentaires
\[
\operatorname{cosupp}X=\bigcap_{Y\in\widetilde X}\deleted{\operatorname{cosupp}}\unsure{\mathbb{C}_M|\mathrm{C}_M}(Y^{\circ})
\]
